% UK Patent Application — Combined Specification
% AMTTP + Adaptive Friction + UDL + BSDT (four inventions, single filing)
% Applicant / Inventor: Odeyemi Olusegun Israel
% Date: 4 March 2026

\documentclass[12pt,a4paper]{article}
\usepackage[margin=2.5cm]{geometry}
\usepackage{amsmath,amssymb}
\usepackage{graphicx}
\usepackage{enumitem}
\usepackage{hyperref}
\usepackage{titlesec}
\usepackage{fancyhdr}
\usepackage{booktabs}
\setlength{\headheight}{14.5pt}

\pagestyle{fancy}
\fancyhf{}
\rhead{Patent Application}
\lhead{AMTTP/AF/UDL/BSDT --- Odeyemi O.I.}
\rfoot{Page \thepage}

\titleformat{\section}{\large\bfseries}{\thesection.}{0.5em}{}
\titleformat{\subsection}{\normalsize\bfseries}{\thesubsection.}{0.5em}{}
\titleformat{\subsubsection}{\normalsize\itshape}{\thesubsubsection.}{0.5em}{}

\begin{document}

\begin{center}
{\LARGE\bfseries PATENT SPECIFICATION}\\[1em]
{\Large State-Adaptive Compliance Enforcement System\\
with Multi-Spectrum Anomaly Detection and\\
Failure-Mode Correction for Decentralised Finance Networks}\\[2em]
{\large Applicant: Odeyemi Olusegun Israel}\\[0.5em]
{\large Inventor:  Odeyemi Olusegun Israel}\\[0.5em]
{\large Date of Filing: 4 March 2026}\\[0.5em]
{\large Application Number: [TO BE ASSIGNED]}\\[0.3em]
{\normalsize \textit{Filed pursuant to the Patents Act 1977}}
\end{center}

\vspace{2em}

% ============================================================
\section{TECHNICAL FIELD}
% ============================================================

The present invention relates to systems and methods for enforcing
anti-money laundering (AML) and regulatory compliance in decentralised
finance (DeFi) transactions using blockchain smart contracts, machine
learning risk scoring, and zero-knowledge cryptographic proofs. More
particularly it relates to a compliance system that performs \\emph{pre-settlement}
transaction flow control in a distributed ledger network by causing smart
contracts to permit, delay (escrow), or reject execution of a transaction
request based on: (i) a machine learning risk score; and (ii) a selected
enforcement profile that deterministically defines numerical risk thresholds
and cryptographic proof requirements to be verified on-chain as a
prerequisite to settlement. In preferred embodiments the risk score is
computed using (a) a multi-spectrum deviation representation that projects
the transaction feature vector through multiple law-domain operators to form
an attribution-rich anomaly tensor, and (b) a failure-mode correction layer
that applies post-hoc correction operators to a frozen deployed model
without retraining. In further embodiments the enforcement profile is
selected in real time by a closed-form state-adaptive controller derived
from a dynamical energy model of the transaction network.

% ============================================================
\section{BACKGROUND}
% ============================================================

\subsection{Technical Problems}

DeFi protocols process financial transactions on blockchain networks
without centralised intermediaries, creating four interacting technical
problems that known art has not resolved in combination.

\subsubsection{Single-Perspective Anomaly Detection Fails to Cover
the Full Space of Fraud Patterns}

Standard anomaly detectors (Isolation Forest, LOF, One-Class SVM,
autoencoders) each produce a single scalar score. They reveal \emph{that}
a transaction is anomalous but not \emph{which mathematical aspect} makes
it so. No single detector spans the full space of mathematical anomaly.
Scores may concentrate in a narrow band, failing to discriminate among
fraud subtypes (the coverage-AUC gap). No prior detector provides
per-law-domain attribution required for downstream failure diagnosis.

\subsubsection{Structural Blind Spots in Deployed Scoring Models}

Supervised models deployed for fraud detection fail to detect certain
fraudulent transactions for structural reasons: mimicry of legitimate
patterns (camouflage), informational sparsity (feature gap), novel
activity scale (activity anomaly), and temporal distributional shift
(temporal novelty). Fewer than 0.1\% of blockchain transactions are
labelled illicit, making retraining impractical. SHAP and LIME explain
individual predictions in input-feature space but do not diagnose
structural failure patterns or provide correction mechanisms without
model retraining.

\subsubsection{Pre-Settlement Enforcement Gap}

Blockchain analytics platforms (Chainalysis, Elliptic) operate
retrospectively, identifying suspicious transactions after settlement
when value transfer is irrecoverable. Protocol-level KYC/AML checks
are bypassed when users transact directly through DeFi smart contracts.
Privacy regulations (UK Data Protection Act 2018, EU GDPR) conflict with
disclosure-intensive compliance models.

\subsubsection{Static Enforcement Parameters}

Existing compliance systems apply fixed risk thresholds and fixed proof
requirements regardless of the network's structural state. This is
provably suboptimal: above a critical manifold in state space any constant
enforcement configuration fails to achieve the state-dependent welfare
optimum. No prior art derives a closed-form optimal enforcement function
mapping a computed energy landscape to numerical thresholds and proof
requirements.

\subsection{Prior Art Limitations}

Known prior art includes: Isolation Forest, LOF, One-Class SVM (single-score,
no attribution); deep autoencoders and ODIN (single latent representation);
SHAP, LIME (feature-wise explanation, no structural correction); conformal
prediction (calibrated sets, no failure diagnosis); blockchain analytics
retrospective monitoring platforms; exchange-level KYC/AML not extending
to protocol transactions; Basel~III countercyclical capital buffer rule
applying a fixed piecewise-linear function to credit-to-GDP gap with no
closed-form state-dependent optimal policy; and agent-based financial
network simulations lacking closed-form optimal policies and provable
stability guarantees. None provides multi-law-domain anomaly tensor
construction, four-component failure-mode correction on a frozen model,
pre-settlement blockchain enforcement, and closed-form state-adaptive
threshold and proof-requirement selection in combination.

% ============================================================
\section{SUMMARY OF THE INVENTION}
% ============================================================

The invention provides a unified state-adaptive compliance enforcement
system for DeFi networks. The shared inventive concept is
\emph{controlling pre-settlement execution of decentralised finance
transactions by selecting, from a predetermined set, an enforcement profile
that deterministically defines numerical risk thresholds and on-chain
cryptographic proof requirements, wherein the selection is computed from
network instability metrics and the resulting policy is enforced by smart
contracts before settlement}. Preferred embodiments further compute the risk
score using a multi-spectrum anomaly tensor and correct structural blind
spots of a frozen deployed model without retraining. The primary technical
contributions are:

\begin{enumerate}[label=(\roman*)]

\item A \textbf{four-layer DeFi compliance architecture} (interface,
compliance orchestration, offline training, infrastructure) enforcing
policy actions before transaction settlement.

\item A \textbf{multi-spectrum transaction anomaly engine} projecting
each transaction feature vector through at least five independent
mathematical law domains---statistical, dynamical, spectral, geometric,
and exponential---to construct a multi-dimensional anomaly tensor
$T_{i,k,d}$, with a formal coverage guarantee; and decomposing that
tensor into magnitude, direction, and novelty components for per-law-domain
attribution.

\item A \textbf{boundary-centred dimension magnifier} rescaling each
tensor dimension by its empirical class-separability, amplifying
well-separated dimensions using a single global hyperparameter set
$(\gamma{=}5.0,\ \lambda{=}10^{-4})$.

\item A \textbf{weak-supervision offline training pipeline} comprising
a composite teacher ensemble (autoencoder--gradient boosting, FATF
rule-based, graph-structural) generating pseudo-labels from unlabelled
blockchain data and an XGBoost student model trained on those labels.

\item A \textbf{failure-mode correction layer} operating on the frozen
deployed student model: it decomposes each transaction into four
near-orthogonal components---Camouflage ($C$), Feature Gap ($G$), Activity
Anomaly ($A$), Temporal Novelty ($T$)---computes a composite blind-spot
score (MFLS) with self-calibrating weights, and applies a post-hoc
correction operator to the student model's predicted probability, producing
a corrected risk score without model retraining.

\item A \textbf{state-adaptive transaction flow control module} computing
one or more network instability metrics from recent on-chain activity and
selecting, from a predetermined set of enforcement profiles, a
configuration deterministically defining: (a) a risk threshold profile
$(\tau_1,\tau_2,\tau_3,\tau_4)$ for the five-tier decision matrix; and
(b) a proof-requirement policy specifying, for at least one tier or
transaction class, one or more cryptographic proofs required on-chain
prior to settlement.

\item A \textbf{gravitational interaction engine} modelling $N$ financial
institutions as particles under a pairwise erf-attraction plus logarithmic
repulsion potential; a \textbf{spectral phase-transition detector}
determining whether the network lies above the critical manifold; a
\textbf{blind-spot energy module} computing a macro-level complementary
energy function from four deviation channels; and a \textbf{closed-form
optimal adaptive policy} $\gamma^*(X)$ derived from the Bellman equation
that drives profile selection.

\item A \textbf{zero-knowledge proof subsystem} configured to verify, on-chain,
one or more privacy-preserving compliance predicates as prerequisites to
settlement according to the proof-requirement policy, wherein the proofs are
verifiable without revealing underlying personal data.

\item \textbf{Smart contract infrastructure} enforcing all policy actions
deterministically before settlement with an immutable audit trail.

\end{enumerate}

% ============================================================
\section{DETAILED DESCRIPTION}
% ============================================================

\subsection{System Architecture Overview}
\label{sec:arch}

The system comprises four hierarchically organised layers.

\textbf{Layer~I --- Interface Layer.}
Client SDKs (TypeScript, Python) and a RESTful API receive transaction
requests, return risk scores, and surface compliance status. A web
dashboard provides visual monitoring.

\textbf{Layer~II --- Compliance Orchestration Layer.}
\begin{itemize}
\item \emph{Multi-spectrum transaction anomaly engine} (\S\ref{sec:udl})
\item \emph{Failure-mode correction layer} (\S\ref{sec:bsdt})
\item \emph{Graph analytics module}: transaction graph construction;
centrality, PageRank, community detection.
\item \emph{Sanctions screening module}: OFAC and EU database checks.
\item \emph{State-adaptive transaction flow control module}
(\S\ref{sec:af}).
\item \emph{Deterministic policy engine}: maps corrected scores to policy
actions via the thresholds of the selected enforcement profile.
\end{itemize}

\textbf{Layer~III --- Offline Training Layer.}
Composite teacher ensemble, pseudo-label generation, student model training,
cross-validation, sigmoid calibration.

\textbf{Layer~IV --- Infrastructure Layer.}
Blockchain smart contracts, privacy-preserving proof circuits, databases,
distributed file storage.

\subsection{Pre-Settlement Smart Contract Enforcement Workflow}
\label{sec:workflow}

In a preferred embodiment the system enforces compliance by controlling
execution of a transaction request at the point of settlement in the
blockchain transaction processing network. A representative workflow is:

\begin{enumerate}[label=(\alph*)]
\item A client application submits a transaction request to a core smart
contract function, the request comprising at least a sender address, a
recipient address, and a transaction value.

\item The compliance orchestration layer computes a risk score from the
feature vector and selects an enforcement profile based on instability
metrics computed over a recent sliding time window of on-chain activity.

\item The selected enforcement profile deterministically defines (i) a set
of numerical thresholds for the compliance decision matrix, and (ii) a
proof-requirement policy specifying which cryptographic proofs must be
verified for one or more threshold tiers or transaction classes.

\item If the proof-requirement policy requires one or more proofs, the
transaction is prevented from settling unless verification succeeds on-chain.
In a preferred embodiment the proof is verified by calling a verifier
contract that records a verification result, a proof identifier, and a
timestamp, and the core contract checks the verification result and that the
verification is fresh within a predetermined validity window and has not been
replayed.

\item Based on the risk score and the selected thresholds, the core contract
either (i) authorises execution, (ii) transfers value into an escrow
contract (time-lock), or (iii) rejects execution by reverting.
\end{enumerate}

This workflow provides a technical effect beyond risk scoring \emph{per se}:
it prevents irreversible settlement of non-compliant transactions by
enforcing deterministic preconditions within the transaction processing
network itself, and it enforces privacy-preserving proof-gating without
requiring disclosure of underlying identity or score values.

% ============================================================
\subsection{Multi-Spectrum Transaction Anomaly Engine (UDL)}
\label{sec:udl}
% ============================================================

\subsubsection{Universal Deviation Principle}

\textbf{Definition (Deviation-Separating Family):}
An operator family $\{\Phi_k\}_{k=1}^{K}$ is
\emph{deviation-separating} if the stacked Jacobian
$J(x) = [D\Phi_1(x)^\top \cdots D\Phi_K(x)^\top]^\top$
has full column rank for all $x$ in the transaction feature domain.

\textbf{Theorem~1 (Coverage Guarantee):}
If $\{\Phi_k\}$ is deviation-separating, no anomalous transaction
can simultaneously match the normal projection under all operators.
This provides a detection guarantee absent from every single-perspective
detector.

\textbf{Theorem~2 (Global Bound):}
On a compact anomaly set a quantitative lower bound on minimum deviation
across all operators can be established.

\textbf{Theorem~3 (Finite-Sample Guarantee):}
Given $n$ normal training samples the coverage guarantee holds with
probability $\geq 1-\delta$ for a sample complexity depending on the
operator family's covering number.

\textbf{Theorem~4 (Information-Optimal Fusion):}
Optimal operator-output weights are the solution to a semidefinite
programme.

\textbf{Theorem~5 (Greedy Operator Selection):}
A greedy algorithm achieves a $(1-1/e)$ approximation to the optimal
operator subset under a submodular information criterion.

\subsubsection{Spectrum Operators}
\label{sec:ops}

Five spectrum operators $\Phi_k$ are instantiated on the 167-dimensional
transaction feature vector:

\paragraph{Statistical Operator ($\Phi_{\mathrm{stat}}$).}
Shannon entropy of the local neighbourhood; KL divergence from the
reference distribution; Hellinger distance; skewness; kurtosis.

\paragraph{Dynamical Operator ($\Phi_{\mathrm{chaos}}$).}
Lyapunov exponent from the temporal transaction context; recurrence rate;
approximate entropy.

\paragraph{Spectral Operator ($\Phi_{\mathrm{freq}}$).}
Power spectral density divergence; low-to-high frequency energy ratio;
spectral entropy; spectral centroid shift.

\paragraph{Geometric Operator ($\Phi_{\mathrm{geom}}$).}
Mahalanobis distance from the reference centroid; cosine similarity to
the reference direction; norm ratio.

\paragraph{Exponential Operator ($\Phi_{\mathrm{exp}}$).}
\begin{equation}
\Phi_{\mathrm{exp}}(x) = \exp\!\left(\alpha\cdot\tfrac{x-\mu}{\sigma}\right)
\end{equation}
amplifying small deviations invisible under linear representations.
$\mu$, $\sigma$ are reference statistics computed on normal transactions.

In an enhanced embodiment up to 14 composable operators are available
(phase-space, topological, kernel, rank, wavelet). An operator correlation
audit confirms at most one pair exceeds the redundancy threshold
$(\rho<0.85)$.

\subsubsection{Anomaly Tensor Construction}

Each transaction feature vector $x_i$ is projected through all $K$
operators to yield a three-dimensional anomaly tensor:
\begin{equation}
T_{i,k,d} \quad\text{for transaction }i,\ \text{operator }k,
\ \text{dimension }d
\end{equation}

\subsubsection{Magnitude--Direction--Novelty Decomposition}

\begin{enumerate}[label=(\alph*)]
\item \textbf{Magnitude} $r_i = \|T_i\|$: total deviation amount.
\item \textbf{Direction} $\hat{D}_i = T_i/\|T_i\|$: unit tensor
indicating which law domains drive the anomaly.
\item \textbf{Novelty} $\theta_i$: angular distance of $\hat{D}_i$ from
reference deviation directions in training data.
\end{enumerate}

\subsubsection{Boundary-Centred Dimension Magnifier}

Each tensor dimension $d$ is rescaled:
\begin{equation}
r'_d = \tanh\!\left(k_d\cdot\tfrac{r_d-b_d}{\sigma_d}\right)
\end{equation}
where $b_d$ is the decision boundary in dimension $d$, $\sigma_d$ is the
local scale, and $k_d = \gamma\cdot\delta_d$ is the per-dimension
steepness. The discriminability score $\delta_d$ is computed from three
complementary measures: Cohen's $d$ (effect size), class-distribution
overlap area, and per-dimension AUC. The global parameters $\gamma{=}5.0$
and QDA regularisation $\lambda{=}10^{-4}$ constitute the single
hyperparameter set used across all deployments.

\subsubsection{Classification and Coverage}

The rescaled tensor is classified by QDA-Magnified (preferred) or a hybrid
auto-selector that evaluates multiple strategies and maximises coverage:
\begin{equation}
\mathrm{Coverage@}k = \frac{|\{x:\mathrm{rank}(x)\leq k,\
x\in\mathrm{anomaly}\}|}{|\mathrm{anomaly}|}
\end{equation}
In the preferred embodiment mean coverage exceeds 91\%. The anomaly
tensor is the primary input to both the failure-mode correction layer
(\S\ref{sec:bsdt}) and the risk scoring engine.

% ============================================================
\subsection{Failure-Mode Correction Layer (BSDT)}
\label{sec:bsdt}
% ============================================================

The failure-mode correction layer receives the anomaly tensor from
\S\ref{sec:udl} and the predicted probability $p(x)$ from the deployed
frozen student model, and produces a corrected probability $p^*(x)$
without modifying the model's parameters.

\subsubsection{Relationship to Spectrum Operators}

The four BSDT components correspond to specialised instantiations of
the UDL spectrum operators applied to the student model's feature space:
Camouflage ($C$) corresponds to the geometric operator applied to the
legitimate-class centroid; Feature Gap ($G$) corresponds to the
statistical operator measuring informational sparsity; Activity Anomaly
($A$) corresponds to the dynamical operator measuring deviation from
known fraud activity; Temporal Novelty ($T$) corresponds to the spectral
and dynamical operators applied over the temporal axis of the training
distribution. At the macro level the same four channel types reappear in
the blind-spot energy module of the systemic risk controller
(\S\ref{sec:af}).

\subsubsection{Camouflage Component}

\begin{equation}
C(x) = 1 - \frac{\|x-\mu_{\mathrm{legit}}\|_2}{d_{\max}}
\end{equation}
$C(x)\approx 1$ indicates that the transaction mimics legitimate patterns,
depressing the student model's suspicion score.

\subsubsection{Feature Gap Component}

\begin{equation}
G(x) = \frac{|\{j:|x_j|<\varepsilon\}|}{d}, \quad \varepsilon=10^{-8}
\end{equation}
High $G(x)$ indicates many uninformative features, depriving the model
of discriminative signal.

\subsubsection{Activity Anomaly Component}

\begin{equation}
A(x) = \sigma\!\left(
\frac{\log(1+|\mathrm{activity}(x)|)-\mu_{\mathrm{ref}}}
{\sigma_{\mathrm{ref}}}\right)
\end{equation}
Measures deviation of transaction activity from the reference
detected-fraud population; captures fraud operating at scales not
represented in training labels.

\subsubsection{Temporal Novelty Component}

\begin{equation}
T(x) = \sigma\!\left(\tfrac{1}{2}(\bar{m}(x)-2)\right)
\end{equation}
$\bar{m}(x)$ is the mean temporal distance to the $k$ nearest training
neighbours, normalised by the temporal span.

\subsubsection{Composite Blind-Spot Score and Self-Calibrating Weights}

\begin{equation}
\mathrm{MFLS}(x) = \sum_{i\in\{C,G,A,T\}} w_i\cdot S_i(x)
\end{equation}
The weights $w_i$ (normalised to sum to 1) are computed by one or more
of: mutual information; Fisher variance ratio; Bayesian online updating;
gradient-based optimisation. The four components are verified
near-orthogonal: mean NMI $\in[0.051,0.064]$ across all pairs (threshold
0.30); all VIF $<1.45$.

\subsubsection{Correction Operators}

Four post-hoc correction operators adjust $p(x)$ based on the MFLS:

\paragraph{Additive.}
\begin{equation}
p^*(x) = p(x) + \lambda\cdot\mathrm{MFLS}(x)\cdot(1-p(x))
\cdot\mathbb{1}[p(x)<\tau]
\end{equation}

\paragraph{SignedLR.}
\begin{equation}
p^*(x) = \mathrm{clip}\!\left(p(x)+\beta_0+\textstyle\sum_i\beta_i S_i(x),0,1\right)
\end{equation}

\paragraph{QuadSurf (preferred).}
\begin{equation}
p^*(x) = \mathrm{clip}\!\left(p(x)+\phi_2([C,G,A,T])^\top\hat{\beta},0,1\right)
\end{equation}
$\phi_2$ is the degree-2 polynomial feature map;
$\hat{\beta}$ are ridge-regression coefficients ($\alpha{=}1.0$).

\paragraph{ExpGate.}
\begin{equation}
p^*(x) = \mathrm{clip}\!\left(p(x)+g(x)\cdot q(x),0,1\right)
\end{equation}
$q(x)$ is the QuadSurf output; $g(x)=\sigma(v^\top[C,G,A,T]+b)$ is a
learned sigmoid gate.

The corrected probability $p^*(x)$ is multiplied by 1000 to yield the
integer risk score used by the decision matrix.

% ============================================================
\subsection{Weak Supervision Offline Training Pipeline}
\label{sec:ws}
% ============================================================

Three independent teacher models generate pseudo-labels from unlabelled
blockchain transaction data:
\begin{enumerate}[label=(\alph*)]
\item \textbf{Autoencoder--Gradient Boosting Teacher} ($s_{\mathrm{AE-XGB}}$):
autoencoder reconstruction error and latent features feed an XGBoost
model.
\item \textbf{FATF Rule-Based Teacher} ($s_{\mathrm{FATF}}$): encodes FATF
red-flag indicators as deterministic scoring rules.
\item \textbf{Graph-Based Teacher} ($s_{\mathrm{Graph}}$): suspicion scores
from proximity to high-risk addresses, temporal clustering, and unusual
connectivity.
\end{enumerate}
Composite teacher score:
\begin{equation}
s_{\mathrm{teacher}}(x) = w_1 s_{\mathrm{AE-XGB}}(x) +
w_2 s_{\mathrm{FATF}}(x) + w_3 s_{\mathrm{Graph}}(x)
\end{equation}
Preferred: $w_1{=}0.4$, $w_2{=}0.3$, $w_3{=}0.3$.
Binary pseudo-labels train an XGBoost student model on 167 engineered
features. Raw output is sigmoid-recalibrated:
\begin{equation}
\hat{p}_{\mathrm{cal}}(x) = \sigma\!\left(k\cdot(s(x)-P_{75})\right)
\end{equation}
where k = 30 and P75 is the 75th percentile of the training
score distribution. The BSDT correction layer (\S\ref{sec:bsdt}) then
produces the final corrected integer risk score.

% ============================================================
\subsection{Compliance Decision Matrix}
\label{sec:matrix}
% ============================================================

The deterministic policy engine maps integer scores $[0,1000]$ to policy
actions according to thresholds drawn from the active enforcement profile:

\begin{center}
\begin{tabular}{lll}
\toprule
\textbf{Score Range} & \textbf{Policy Action} & \textbf{Smart Contract Effect}\\
\midrule
$[0,\tau_1)$ & APPROVE & Immediate execution \\
$[\tau_1,\tau_2)$ & APPROVE + MONITOR & Execution; monitoring flag \\
$[\tau_2,\tau_3)$ & REVIEW & Queued for manual review \\
$[\tau_3,\tau_4)$ & ESCROW & Time-locked escrow \\
$[\tau_4,1000]$ & BLOCK & On-chain rejection \\
\bottomrule
\end{tabular}
\end{center}

Under the stable profile: $(\tau_1,\tau_2,\tau_3,\tau_4) =
(200,400,600,800)$. The state-adaptive module (\S\ref{sec:af}) selects
profiles that adjust these thresholds and impose additional proof
requirements.

\subsubsection{Example Enforcement Profiles}

In a preferred embodiment the predetermined set of enforcement profiles
includes at least a \textit{stable} profile, a \textit{stressed} profile, and
a \textit{critical} profile. Each profile deterministically defines both a
threshold tuple $(\tau_1,\tau_2,\tau_3,\tau_4)$ and a proof-requirement
policy.

\begin{center}
\begin{tabular}{llll}
\toprule
\textbf{Profile} & $(\tau_1,\tau_2,\tau_3,\tau_4)$ & \textbf{Proof-gated tiers} & \textbf{Required proofs}\\
\midrule
Stable & (200, 400, 600, 800) & $[\tau_3,1000]$ & non-membership predicate \\
Stressed & (250, 450, 650, 850) & $[\tau_2,1000]$ & non-membership + range predicate \\
Critical & (300, 500, 700, 900) & $[\tau_1,1000]$ & non-membership + range + credential predicate \\
\bottomrule
\end{tabular}
\end{center}

The example profiles are illustrative and may be modified for a given
deployment. For a fixed system state, the selected profile and the resulting
on-chain decisions are deterministic and reproducible.

% ============================================================
\subsection{State-Adaptive Transaction Flow Control Module}
\label{sec:af}
% ============================================================

\subsubsection{Instability Metrics}

Network instability metrics are computed from:
\begin{enumerate}[label=(\alph*)]
\item spectral radius or centrality concentration of the transaction graph
over a sliding time window;
\item a temporal burst indicator from the transaction count process; and/or
\item a shift of the corrected risk score distribution relative to its
reference distribution.
\end{enumerate}

\subsubsection{Enforcement Profile Selection}

The module selects one profile from the discrete set
$\{\textit{stable},\, \textit{stressed},\, \textit{critical}\}$.
Each profile deterministically defines:
\begin{enumerate}[label=(\alph*)]
\item a risk threshold profile $(\tau_1,\tau_2,\tau_3,\tau_4)$; and
\item a proof-requirement policy specifying which privacy-preserving
compliance predicate proofs must be verified on-chain for at least one tier
or transaction class, the compliance predicates comprising at least one of:
(i) a set non-membership predicate, (ii) a numerical range predicate, or
(iii) a credential validity predicate.
\end{enumerate}

\subsubsection{Gravitational Interaction Engine}

When multi-institution financial metrics are available, the module models
$N$ institutions as particles in $\mathbb{R}^d$:
\begin{equation}
\dot{X} = -\alpha(X-\mu\mathbf{1}^\top) - \nabla_X\Phi_{\mathrm{pair}}(X)
- \sum_k \beta_k D\Phi_k(X)^\top D_k(X)
\end{equation}

\subsubsection{Pairwise Interaction Potential}

\begin{equation}
\phi(r) = \gamma\cdot\frac{\sigma\sqrt{\pi}}{2}\,\mathrm{erf}(r/\sigma)
- \gamma\lambda_{\mathrm{rep}}\cdot\log(r+\varepsilon)
\end{equation}
Preferred parameters: $\alpha{=}0.10$, $\gamma{=}1.00$,
$\sigma{=}1.00$, $\lambda_{\mathrm{rep}}{=}0.10$, $\varepsilon{=}10^{-5}$.

\subsubsection{Macro-Level Blind-Spot Energy Module}

A macro-level energy function is computed from four deviation channels
that mirror, at the network level, the BSDT components at the transaction
level:
\begin{enumerate}[label=(\alph*)]
\item enstrophy anomaly $\delta_C$: rotational energy deviation (analogous
to Camouflage);
\item spectral gap $\delta_G$: structural fragility indicator (analogous
to Feature Gap);
\item alignment anomaly $\delta_A$: co-movement deviation (analogous to
Activity Anomaly); and
\item temporal novelty $\delta_T$: drift from reference state patterns
(analogous to Temporal Novelty).
\end{enumerate}

\begin{equation}
E_{\mathrm{BS}}(X) =
\|(X-\mu_0)^\top\Sigma_0^{-1}(X-\mu_0)\|_F
\end{equation}
$\mu_0$ and $\Sigma_0$ are reference mean and covariance estimated by
Ledoit--Wolf Oracle shrinkage.

\subsubsection{Critical Manifold}

\textbf{Theorem (Spectral Phase Transition):}
\begin{equation}
\mathcal{C} = \{X:\lambda_{\max}(\rho(X)\cdot\ell(X)\cdot W)=1\}
\end{equation}
Above $\mathcal{C}$ state-adaptive friction is both necessary and
sufficient for welfare optimality.

\subsubsection{Closed-Form Optimal Adaptive Policy}

\textbf{Theorem (Optimal Friction):} The Bellman-optimal friction
coefficient is:
\begin{equation}
\gamma^*(X) \approx
\frac{\beta\|\nabla E\|^2}{2\kappa+\beta\|\nabla E\|^2}
\cdot\frac{E(X)}{E(X)+\theta}
\end{equation}
$\gamma^*\to 0$ as $E\to 0$ (minimum friction when stable);
$\gamma^*\to 1$ as $E\to\infty$ (maximum friction under severe
instability). The adaptive damping factor $E/(E+\theta)$ is the same
functional form as the BSDT damping function $\gamma^*(E)=E/(E+\theta)$
applied at the transaction level by the correction operators, linking
both levels of the system through a single analytical form.

Profile selection: below $\gamma_1$ --- \textit{stable};
between $\gamma_1$ and $\gamma_2$ --- \textit{stressed};
above $\gamma_2$ --- \textit{critical}.

\subsubsection{Navier--Stokes Analogy}

Institution position changes $\dot{X}$ are interpreted as a velocity
field decomposed into strain $S$ and vorticity $\Omega$. Enstrophy
$\mathcal{E}=\frac{1}{2}\|\omega\|_F^2$ quantifies rotational energy.
The Beale--Kato--Majda criterion adapted to financial networks serves
as a singularity indicator. The adaptive friction coefficient is applied
as an adaptive viscosity $\nu=\nu_0(1+\gamma^*(E_{\mathrm{BS}}))$.
Minsky regime classification (hedge, speculative, Ponzi) emerges from
the cosine alignment between the gravitational and blind-spot gradients.

% ============================================================
\subsection{Zero-Knowledge Proof Subsystem (zkNAF)}
\label{sec:zknaf}
% ============================================================

In preferred embodiments, the system includes one or more zero-knowledge proof
circuits implementing privacy-preserving compliance predicates. In one
embodiment the proof system is a succinct non-interactive argument of
knowledge (SNARK), for example Groth16.

\textbf{Example: Set Non-Membership Predicate.}
The prover shows an address is absent from a cryptographic set commitment
without revealing the address.

\textbf{Example: Numerical Range Predicate.}
The prover demonstrates $s_{\min}\leq s(x)\leq s_{\max}$ without
revealing the exact corrected score.

\textbf{Example: Credential Validity Predicate.}
The prover demonstrates satisfaction of an identity or credential status with
an accredited provider without revealing identity data.

Each proof may carry a predetermined validity interval. An on-chain verifier
contract performs verification and records the result together with a
timestamp and a proof identifier.

% ============================================================
\subsection{Smart Contract Architecture}
\label{sec:sc}
% ============================================================

\begin{itemize}
\item \textbf{Core Contract:} on-chain decision matrix; immutable audit
record.
\item \textbf{Escrow Contract:} time-locked escrow; release and dispute
escalation.
\item \textbf{ZK Verifier Contract:} verification of selected zero-knowledge
proofs; binary result stored with timestamp and proof identifier.
\item \textbf{Cross-Chain Bridge Contract:} requires compliance
verification on both source and destination chains.
\item \textbf{Dispute Resolution Contract:} structured escalation for
escrowed disputes.
\end{itemize}

All enforcement decisions execute before the corresponding transaction
settles on the blockchain network.

% ============================================================
\subsection{Performance Characteristics}
\label{sec:perf}
% ============================================================

\textbf{DeFi Compliance System} (2.64 million Ethereum transactions):
\begin{center}
\begin{tabular}{ll}
\toprule
\textbf{Metric} & \textbf{Value} \\
\midrule
Meta-Ensemble ROC-AUC (test) & 1.0000 \\
Meta-Ensemble ROC-AUC (5-fold CV) & $0.9965\pm0.0049$ \\
Meta-Ensemble PR-AUC & 0.9999 \\
Transaction XGBoost AUC & 0.9878 \quad F1: 0.901 \\
Cross-domain mean AUC (9 datasets) & 0.9962 \\
Escrow contract gas cost & 287,890 gas \\
\bottomrule
\end{tabular}
\end{center}

\textbf{UDL Anomaly Engine} (single hyperparameter set, $\gamma{=}5.0$):
\begin{center}
\begin{tabular}{lcc}
\toprule
\textbf{Dataset} & \textbf{QDA-Magnified AUC} & \textbf{Hybrid AUC}\\
\midrule
Synthetic & 1.000 & 0.999 \\
Mimic (adversarial) & 1.000 & 0.999 \\
Mammography & 0.947 & 0.911 \\
Shuttle & 0.989 & 0.991 \\
Pendigits & 0.777 & 0.960 \\
Mean & 0.943 & 0.972 \\
\bottomrule
\end{tabular}
\end{center}

\textbf{BSDT Correction Layer} (792 variant--model--dataset combinations):
\begin{center}
\begin{tabular}{ll}
\toprule
\textbf{Metric} & \textbf{Value} \\
\midrule
Best variant (QuadSurf) mean F1 & 0.787 ($+43\%$ relative) \\
False positive reduction & 88.1\% (44,740 $\to$ 5,345) \\
Maximum recall recovery & $+84.4$ percentage points \\
Shuttle rescue F1 & $0.134\to 0.968$ \\
Cross-domain credit card MFLS AUC & 0.885 \\
Cross-domain aerospace MFLS AUC & 0.982 \\
Cross-domain medical imaging MFLS AUC & 0.916 \\
\bottomrule
\end{tabular}
\end{center}

\textbf{Adaptive Friction Controller} (FDIC SDI data, 140 quarters):
\begin{center}
\begin{tabular}{lccc}
\toprule
\textbf{Configuration} & $N$ & \textbf{AUROC} & \textbf{GFC Lead}\\
\midrule
Institution-type (7 classes) & 7 & $\geq0.70$ & 6 quarters \\
Bank-level (individual) & 30 & 0.805 & 6 quarters \\
G-SIB cross-border & 20 & 0.781 & 6 quarters \\
\bottomrule
\end{tabular}
\end{center}
Super-criticality: 68.6\% of sample above $\mathcal{C}$.
GFC welfare loss of inaction: 19.57 pp consumption-equivalent.

% ============================================================
\clearpage
\section{CLAIMS}
% ============================================================

\noindent\textit{The following claims define the scope of protection
sought. Features of any dependent claim may be combined with any
independent claim or with other dependent claims.}
\vspace{0.8em}

\begin{enumerate}[label=\textbf{\arabic*.}]

%──────────────────────────────────────────────────────────
% Independent Claim 1 — DeFi Compliance System (Technical Transaction Control)
%──────────────────────────────────────────────────────────
\item A computer-implemented transaction compliance system for a decentralised
finance transaction processing network, the system comprising:
\begin{enumerate}[label=(\alph*)]

\item an interface layer configured to receive, from a client application,
a transaction request comprising at least a sender address, a recipient
address, and a transaction value;

\item a risk scoring engine configured to compute a numerical risk score
for the transaction request by:
\begin{enumerate}[label=(\roman*)]
  \item constructing a feature vector comprising address-level,
  transaction-level, and graph-level features;
  \item projecting the feature vector through a plurality of spectrum
  operators each computing deviation measurements under a different
  mathematical law domain to construct a multi-dimensional anomaly tensor;
  \item applying a trained machine learning model to the anomaly tensor to
  compute the numerical risk score;
\end{enumerate}

\item a state-adaptive transaction flow control module configured to
compute one or more network instability metrics from recent transaction
activity and to compute a state-dependent friction coefficient
$\gamma^*(X)$ based on the one or more network instability metrics, and to
select, from a predetermined set of enforcement profiles, an enforcement
configuration based on the computed friction coefficient $\gamma^*(X)$,
the enforcement configuration comprising:
\begin{enumerate}[label=(\roman*)]
  \item a risk threshold profile defining numerical thresholds for a
  deterministic compliance decision matrix; and
  \item a proof-requirement policy defining, for at least one threshold tier
  or transaction class, one or more cryptographic proofs required to be
  verified on-chain as a prerequisite to settlement;
\end{enumerate}

\item a deterministic policy engine configured to map the risk score to a
policy action according to the selected risk threshold profile and to
enforce the proof-requirement policy by causing a settlement-preventing
action when required proofs are absent, expired, or fail on-chain
verification, wherein the settlement-preventing action comprises at least
one of: causing a core smart contract function to revert execution of a
transaction call, or preventing release of value from an escrow smart
contract; and

\item smart contract infrastructure deployed on a blockchain network and
configured to implement, before settlement, the policy action determined by
the deterministic policy engine, including a core contract configured to
reject or authorise execution, an escrow contract configured to time-lock
value transfers for an escrow policy action, and a zero-knowledge proof
verifier contract configured to verify said cryptographic proofs,
wherein the core contract is further configured to require, for a
transaction class or threshold tier specified by the proof-requirement
policy, a successful verification result from the zero-knowledge proof
verifier contract prior to permitting execution;
\end{enumerate}

\noindent wherein the smart contract infrastructure prevents settlement of
the transaction request unless the policy action permits execution and, when
required by the proof-requirement policy, the cryptographic proofs are
verified on-chain.

%──────────────────────────────────────────────────────────
% Independent Claim 2 — DeFi Compliance Method (Pre-Settlement Enforcement)
%──────────────────────────────────────────────────────────
\item A computer-implemented method of enforcing compliance in decentralised
finance transactions on a blockchain network, the method comprising:
\begin{enumerate}[label=(\alph*)]

\item receiving a transaction request comprising at least a sender address,
a recipient address, and a transaction value;

\item computing a numerical risk score by:
\begin{enumerate}[label=(\roman*)]
  \item constructing a feature vector comprising address-level,
  transaction-level, and graph-level features;
  \item projecting the feature vector through a plurality of spectrum
  operators each computing deviation measurements under a different
  mathematical law domain to construct a multi-dimensional anomaly tensor;
  \item applying a trained machine learning model to the anomaly tensor to
  compute the numerical risk score;
\end{enumerate}

\item computing one or more network instability metrics from recent
transaction activity, computing a state-dependent friction coefficient
$\gamma^*(X)$ based on the one or more network instability metrics, and
selecting an enforcement profile from a predetermined set based on the
computed friction coefficient $\gamma^*(X)$, the selected profile
deterministically defining:
\begin{enumerate}[label=(\roman*)]
  \item a risk threshold profile for a deterministic decision matrix; and
  \item a proof-requirement policy specifying one or more cryptographic
  proofs to be verified on-chain as a prerequisite to settlement for at
  least one threshold tier or transaction class;
\end{enumerate}

\item mapping the risk score to a policy action using the decision matrix
with thresholds from the selected risk threshold profile;

\item before settlement, enforcing the proof-requirement policy by
verifying the specified cryptographic proofs on-chain using a zero-knowledge
proof verifier smart contract and preventing execution of the transaction
when required proofs are absent, expired, or fail verification, wherein
preventing execution comprises reverting a core smart contract function call
or preventing release of value from an escrow smart contract; and

\item executing a corresponding smart contract function on a blockchain
network to implement the policy action, including at least one of:
authorising execution, transferring value into escrow, maintaining value in
escrow until a release condition is satisfied, or rejecting execution by
reversion; and recording the outcome on-chain as an immutable audit record.
\end{enumerate}

%──────────────────────────────────────────────────────────
% Dependent Claims
%──────────────────────────────────────────────────────────
\item A system according to claim 1 or a method according to claim 2,
wherein computing the risk score comprises projecting the feature vector
through a plurality of spectrum operators each computing deviation
measurements under a different mathematical law domain to construct a
multi-dimensional anomaly tensor, and using the anomaly tensor as an input
to the trained machine learning model.

\item A system or method according to claim 3, wherein the plurality of
spectrum operators comprises at least: a statistical operator measuring
distributional deviations, a dynamical operator measuring temporal dynamical
deviations, a spectral operator measuring frequency-domain deviations, a
geometric operator measuring geometric relationship deviations, and an
exponential operator measuring deviations under exponential amplification.

\item A system or method according to claim 3 or claim 4, wherein the
anomaly tensor is decomposed into magnitude, direction, and novelty
components, and the direction and novelty components provide per-law-domain
attribution identifying which law domains contribute to a risk assessment.

\item A system or method according to any of claims 3 to 5, wherein the
spectrum operators collectively satisfy a deviation-separating condition
such that the stacked Jacobian of the operators has full column rank over
the transaction feature domain, providing a coverage guarantee that no
anomalous transaction can simultaneously match normal data across all
operators.

\item A system or method according to any of claims 1 to 6,
wherein the boundary-centred dimension magnifier rescales each anomaly
tensor dimension by:
\[
r'_d = \tanh\!\left(k_d\cdot\tfrac{r_d-b_d}{\sigma_d}\right)
\]
where $k_d=\gamma\cdot\delta_d$, $\delta_d$ is a per-dimension
discriminability score computed from Cohen's $d$, class-distribution
overlap, and per-dimension AUC, and $\gamma$ is a single global steepness
parameter.

\item A system according to claim 1 or a method according to claim 2,
wherein computing the risk score further comprises computing, from an anomaly
tensor and a predicted probability output by a frozen deployed student model,
four near-orthogonal failure-mode components comprising camouflage, feature
gap, activity anomaly and temporal novelty, aggregating them into a
composite blind-spot score, and applying a correction operator to the
predicted probability to produce a corrected risk score without retraining
or updating parameters of the trained machine learning model.

\item A system or method according to claim 8,
wherein the correction operator is a QuadSurf correction:
\[
p^*(x) = \mathrm{clip}\!\left(p(x)+\phi_2([C,G,A,T])^\top\hat{\beta},0,1\right)
\]
where $\phi_2$ is the degree-2 polynomial feature map and $\hat{\beta}$
are ridge-regression coefficients.

\item A system or method according to claim 8,
wherein the correction operator is an exponential gating correction
comprising a QuadSurf output modulated by a learned sigmoid gate
$g(x)=\sigma(v^\top[C,G,A,T]+b)$.

\item A system or method according to any of claims 8 to 10, wherein the
four failure-mode components are aggregated as:
\[
\mathrm{MFLS}(x)=\textstyle\sum_{i\in\{C,G,A,T\}} w_i\cdot S_i(x)
\]
and the weights $w_i$ are automatically computed using one or more of
mutual information, Fisher variance ratio, Bayesian online updating, or
gradient-based optimisation.

\item A system or method according to claim 1 or claim 2, wherein the
enforcement profile is selected from a discrete set and the selection
is determined by whether the computed friction coefficient $\gamma^*(X)$
lies between predetermined lower and upper thresholds.

\item A system or method according to claim 1 or claim 2, wherein the
proof-requirement policy comprises, for at least one enforcement profile,
requiring at least one privacy-preserving cryptographic proof establishing
satisfaction of a compliance predicate, wherein the compliance predicate
comprises at least one of: (i) a set non-membership predicate, (ii) a
numerical range predicate, or (iii) a credential validity predicate.

\item A system or method according to claim 13, wherein the
verifier smart contract records, for a successful proof verification, at
least a proof identifier and a verification timestamp, and the core smart
contract prevents settlement when the proof identifier has been previously
used or when the verification timestamp is older than a predetermined
validity threshold.

\item A system according to claim 1 or a method according to claim 2,
wherein the proof-requirement policy deterministically defines one or more
proof types as a function of (i) the computed state-dependent friction
coefficient $\gamma^*(X)$ and (ii) a per-law-domain attribution derived from
the anomaly tensor, such that at least one of: (a) $\gamma^*(X)$ exceeding a
friction threshold or (b) an attribution component exceeding an attribution
threshold causes the deterministic policy engine to require at least one of:
a set non-membership predicate proof, a numerical range predicate proof, or a
credential validity predicate proof.

\item A system or method according to claim 1 or claim 2, wherein the
state-dependent friction coefficient is:
\[
\gamma^*(X)\approx
\frac{\beta\|\nabla E\|^2}{2\kappa+\beta\|\nabla E\|^2}
\cdot\frac{E(X)}{E(X)+\theta}
\]
where $E(X)$ is the combined systemic energy, $\kappa$ is an economic
cost parameter, $\beta$ is a systemic damage coefficient, and
$\theta=\kappa/\beta$; and wherein the adaptive damping factor
$E(X)/(E(X)+\theta)$ is the same functional form applied by the
failure-mode correction module at the transaction level.

\item A system or method according to claim 16, wherein a state-adaptive
transaction flow control module comprises a gravitational interaction engine
configured to model financial institutions as particles interacting through a
pairwise potential, a spectral monitoring module configured to determine a
phase-transition condition, and an adaptive policy module configured to
compute the friction coefficient.

\item A system or method according to claim 1 or claim 2, wherein selecting
the enforcement profile further deterministically defines at least one of:
a transaction rate limit, a queue delay, an escrow duration, an escrow
release condition window, or a maximum transaction value limit.

\item A system according to claim 1, further comprising a cross-chain
bridge contract requiring compliance verification on both a source chain
and a destination chain before releasing bridged funds.

\item A system according to claim 1 or a method according to claim 2,
wherein computing the one or more network instability metrics comprises
constructing, over a sliding time window of on-chain activity, a
transaction graph and computing at least one of: (i) a spectral radius of an
adjacency or interaction matrix derived from the transaction graph, (ii) a
centrality concentration metric, or (iii) a temporal burst metric from a
transaction count process.

\item A system according to claim 1 or a method according to claim 2,
wherein computing the one or more network instability metrics comprises
modelling, over the sliding time window, aggregated transaction flow as a
fluid-like flow field and computing at least one of: a divergence metric, a
vorticity metric, or an effective viscosity parameter indicative of
liquidity friction, and selecting the enforcement profile based on the at
least one computed metric.

\item A system according to claim 1 or a method according to claim 2,
wherein computing the one or more network instability metrics comprises
classifying at least one entity, address cluster, or counterparty into a
Minsky regime comprising at least one of: hedge, speculative, or Ponzi,
based on refinancing dependence inferred from on-chain behaviour, and
selecting the enforcement profile based on the classification.

\item A system according to claim 1 or a method according to claim 2,
wherein the predetermined set of enforcement profiles is stored in an
on-chain data structure of the core contract as a mapping from a profile
identifier to (i) the risk threshold profile and (ii) the proof-requirement
policy, and wherein selecting the enforcement profile comprises selecting a
said profile identifier.

\item A system according to claim 1 or a method according to claim 2,
wherein enforcing the proof-requirement policy comprises: (i) calling a
zero-knowledge proof verifier contract that stores a verification result, a
proof identifier, and a verification timestamp, and (ii) checking in the core
contract that the stored verification result indicates success, that the
proof identifier has not been previously used, and that the verification
timestamp is within a predetermined validity window prior to permitting
execution or release from escrow.

\item A system according to claim 1, wherein the core contract emits an
on-chain event log comprising at least: a transaction identifier, a selected
profile identifier, a policy action identifier, and a proof-verification
status, thereby providing an immutable audit trail of pre-settlement
enforcement.

\end{enumerate}

% ============================================================
\newpage
\section*{ABSTRACT}
% ============================================================

A computer-implemented system controls pre-settlement execution of DeFi
transactions on a blockchain network. A risk scoring engine constructs an
anomaly tensor by applying multiple spectrum operators to transaction
features and applies a trained model to compute a risk score. A
state-adaptive module computes network instability metrics and a
state-dependent friction coefficient $\gamma^*(X)$ and selects an
enforcement profile that sets risk thresholds and a proof-requirement
policy. Smart contracts enforce the profile before settlement by reverting
execution or time-locking value in escrow unless required zero-knowledge
proofs are verified on-chain. A verifier contract stores verification
results with timestamps and proof identifiers and the core contract checks
freshness within a predetermined validity window and prevents replay.
On-chain events record the selected profile and
outcome as an immutable audit trail.

\end{document}
